\begin{tikzpicture}[
    dim_line/.style={{Stealth[scale=0.8]}-{Stealth[scale=0.8]}, thin},
    tube_wall/.style={line width=1.5pt, draw=black, line join=round}
]

    % 1. Definimos la ruta del contorno EXACTO para usarla como máscara (clip)
    \def\mascara{
        (-3.5, 1.5) [rounded corners=3mm] -- (-1, 1.5) [rounded corners=0mm] --
        (-1, 3.3) -- (0.3, 3.3) [rounded corners=4mm] --
        (0.3, 0) [rounded corners=0mm]
        arc[start angle=180, end angle=360, radius=1.2] --
        (2.7, 6.5) -- (2.3, 6.5) --
        (2.3, 0)
        arc[start angle=360, end angle=180, radius=0.8] --
        (0.7, 3.7) [rounded corners=0mm] --
        (-1, 3.7) [rounded corners=3mm] --
        (-1, 4.5) [rounded corners=3mm] -- (-3.5, 4.5) -- cycle
    }

    % 2. Pintamos los fluidos usando la máscara para que no se salgan
    \begin{scope}
        \clip \mascara;
        
        % Fluido de la cámara A (celeste claro)
        \fill[cyan!15] (-4, 0.5) rectangle (1, 5);
        
        % Fluido manométrico (gris oscuro)
        \fill[gray!50] (0, 0.5) rectangle (1.5, -2); % Rama izquierda
        \fill[gray!50] (1.5, -2) rectangle (3, 5);   % Rama derecha y base en U
    \end{scope}

    % 3. Dibujamos las paredes del tubo (dejando abierta la parte superior D)
    \draw[tube_wall]
        (2.3, 6.5) --
        (2.3, 0)
        arc[start angle=360, end angle=180, radius=0.8] --
        (0.7, 3.7) [rounded corners=3mm] --
        (-1, 3.7) [rounded corners=3mm] --
        (-1, 4.5) [rounded corners=3mm] -- 
        (-3.5, 4.5) [rounded corners=3mm] -- 
        (-3.5, 1.5) [rounded corners=3mm] -- 
        (-1, 1.5) [rounded corners=0mm] --
        (-1, 3.3) -- (0.3, 3.3) [rounded corners=4mm] --
        (0.3, 0) [rounded corners=0mm]
        arc[start angle=180, end angle=360, radius=1.2] --
        (2.7, 6.5);

    % 4. Meniscos y superficies en los tubos
    \draw[thick] (0.3, 0.5) -- (0.7, 0.5); % Interfaz recta en B
    \draw[thick] (2.3, 5) to[bend right=15] (2.7, 5); % Menisco en D

    % 5. Líneas punteadas de referencia
    \draw[dashed] (-0.5, 0.5) -- (4, 0.5) node[right, font=\small] {N} node[left=4.5cm, font=\small] {M};
    \draw[dashed] (-1, 3.5) -- (2, 3.5); % Nivel de A
    \draw[dashed] (2.7, 5) -- (4, 5); % Nivel de D

    % 6. Cotas de altura
    \draw[dim_line] (1.5, 0.5) -- (1.5, 3.5) node[midway, fill=white, inner sep=2pt] {$h_1$};
    \draw[dim_line] (3.5, 0.5) -- (3.5, 5) node[midway, fill=white, inner sep=2pt] {$h$};

    % 7. Textos y etiquetas clave
    \node at (-2.8, 2) {\Large $\gamma_L$};
    \node at (-1.3, 3.5) {A}; 
    \node[above left] at (0.2, 0.5) {B};
    \node[above left] at (2.2, 0.5) {C};
    \node[left] at (2.2, 5) {D};
    
    % Flecha de Presión atmosférica
    \draw[-Stealth, thick] (2.5, 6.2) -- (2.5, 5.2) node[above=0.2cm, pos=0] {$p_a$};
    
    % Etiqueta del fluido manométrico
    \node[below] at (1.5, -1.3) {$\gamma'_m$};
    \draw[-Stealth] (1.5, -1.3) -- (1.5, -1.2);

\end{tikzpicture}
